% 本章仅引用既有推导与冻结实验；不把模块消融称作连续参数扫描。
\section{模型评价与改进}\label{sec:evaluation}
四问形成了“位置约束、检测点选择、联合航路、无残留认证”的递进体系。评价时，既考察定位与退出判据是否可靠，也考察取得这些证据需要付出多少时间；以下结论依据前文推导及统一离线实验。

\subsection{参数敏感性与稳健性}
{\heiti 误差界影响定位尺度。}\songti 固定检测点与示向度读数，记误差半宽为 $\varepsilon$ 时的角楔交集为 $\mathcal R(\varepsilon)$。由角楔随误差界增大而扩张，有
\begin{equation}\label{eq:evaluation-error}
 \varepsilon_1\le\varepsilon_2
 \ \Longrightarrow\ 
 \mathcal R(\varepsilon_1)\subseteq\mathcal R(\varepsilon_2),\qquad
 D(\varepsilon_1)\le D(\varepsilon_2).
\end{equation}
直径关系针对非空有界交集；误差界继续增大还可能使区域无界。因此，低估误差可能破坏真值包含性，高估误差则可能增加补测成本。问题二的小角近似也表明，远离共线且固定几何与先验时，定位尺度随误差界近似同比增长；近共线布点会放大误差，不能将这一近似推广到全部候选位置。

{\heiti 网格精度影响认证的保守程度。}\songti 问题三由\cref{eq:q3-cert-radius} 得到有效认证半径 $1000-h/\sqrt2-\delta_r$；问题四由\cref{eq:q4-cell,eq:q4-arc} 可见，增大格宽 $h$ 会收紧距离条件并缩短可排除朝向区间。细网格更易取得证书，但增加运算量。问题四的 21 站布局经未证明方格逐级细分，在 2.5 m 通过核验，说明自适应精度可减小几何离散的保守损失；它不意味着该分辨率对所有布局均足够。

{\heiti 改进机制的收益具有场景依赖性。}\songti 单因素消融实验中，问题三取消测站精修后平均增时 \QThreeAblationPolishIncrease\%，问题四取消后验中心估计后增时 \QFourPosteriorChange\%，对应配对区间均为正；但取消问题三的清除后扫描、问题四的信息代价后，时间变化区间均跨零（\cref{fig:q3-cross}(b)、\cref{fig:q4-robustness}(a)）。这支持保留有效机制，也表明不能把每项改进都视为稳定提速来源；上述实验比较了各机制是否采用，尚不能确定连续参数的最优取值。

\subsection{模型优点}
{\heiti 一是定位与完成判据具有可验证依据。}\songti 问题一直接利用误差界构造可行域，揭示直径圆可能不能包住定位区域；后续采用保守外包及包围半径判断清除条件。问题三、四进一步以实际负观测构造连续覆盖证书，使“已知目标清完”与“区域内无残留”得到明确区分。

{\heiti 二是将定位精度与完整任务成本统一考虑。}\songti 问题二以期望直径选择补测位置，问题三、四将测站、清除目标及检测调度纳入航路。在随机场景中，问题三的联合策略 \textrm{joint} 较初代策略 \textrm{field}、问题四的紧凑策略 \textrm{compact} 较前一代策略 \textrm{joint} 的完整时间平均缩短 \QThreeSaving\% 和 \QFourSaving\%；两项比较各自使用相同场景配对，支持联合规划降低移动与检测开销的作用。

{\heiti 三是数值结果具有可复核的实验依据。}\songti 几何计算采用独立方法交叉核验，策略比较保留失败场景，并报告完整退出时间、清除比例及配对区间。问题三的最终策略在 120 组场景全部清除目标，问题四采用连续证书约束的最终策略在 60 组场景清除全部 \QFourSourceTotal\ 个源并完成认证；逐动作记录和固定随机种子使结果能够复核。

\subsection{模型局限}
{\heiti 首先，可靠性依赖观测模型成立。}\songti 保守定位及无源证书要求源静止、接收半径下界有效、误差有界且接收区内无漏报。遮挡、异常测向或源移动均可能破坏这些条件。证书成立能保证无残留，但尚未证明算法在任意合法场景和有限预算内都能取得证书；离线样本全清也不等于未来零失败风险。

{\heiti 其次，估价模型存在近似与先验偏差。}\songti 问题二的面积先验、小角近似及接收可达性处理限制其适用范围，解析值与数值值在已比较位置吻合，不能替代全域误差界。问题四将单方位位置近似为中心射线，并采用固定全向比例和参考状态估价；源分布变化时，规划中心和检测收益可能失准。问题二的 60 m 阈值是期望性能标准，也不保证每次定位直径均低于该值。

{\heiti 最后，求解效率与最优性仍有局限。}\songti 贪心选点、有限候选及局部航路交换不能保证全局最优，21 站也未证明为最少站数。更强的搜索和更细的认证会增加计算开销，光学清除网格在定位区域较大时可能需要较多次尝试。本地运行时间未计入真实接口通信、网络抖动等影响；问题三、四各三次正式测试提供了具体案例的结果，但样本有限，仍不足以推断任意场景下的运行表现。

\subsection{改进方向}
改进应分别对应上述局限。其一，在位置估价中比较不同空间分布、接收半径及全向比例先验，并对近共线、不同交会角分支和接收截断重新核验几何近似；以信息充分的补测位置降低先验依赖，同时保留保守外包与连续证书。

其二，以同一完整时间目标协调测站选择、补测时机和清除顺序，通过小规模精确求解建立比较基准，再评估局部搜索的改进空间；对覆盖边界采用自适应细分，并设置计算预算，平衡虚拟任务耗时与实际运算时间。

其三，在独立场景中扩大误差场、定向比例和边缘源分布的验证范围，报告超时、漏清及接口失败；若引入漏报或移动源，应重建相应的观测与置信判据，经校准后再用于退出决策，不能直接沿用现有确定性无源证书。\label{sec:evaluation-end}
